% === ABSTRACT & KEYWORDS =====================================================
\begin{abstract}
On-track system identification learns tire models from ordinary racing laps
instead of dedicated manoeuvres. We study one such residual-learning pipeline
at full vehicle scale, in a simulator publishing per-wheel forces as ground
truth. Transplanted unchanged from a 1:10 platform it converges on a tire with
less than half the vehicle's grip ceiling and most Pacejka coefficients railed.
The cause is a loss of excitation created by the pipeline's own virtual-data
step, which no change of solver or tuning removes. The friction the
synthetic rollout can demand,
$\mu_{\mathrm{reach}}=v^{2}\delta_{\mathrm{sweep}}/(gL)$, is a property of the
rollout configuration and not of the tire, as a grid on a known tire confirms.
Two further defects share that source: the fit's objective depends on the
coefficients only through the product $BCD$, so the peak factor is not
recoverable from the rollout at any excitation, and the loop feeding each fit
its own previous answer is not a contraction. We correct all three (rollout
speed from the measured lap, sweep bounded to the observed steering, peak
factor measured off the trace and pinned, update damped into a relaxed
fixed-point iteration) and gate the controller interface on derived physical
quantities rather than on coefficient bounds. Over three timed laps on a surface
\SI{33}{\percent} below the one its prior was measured on, no accepted
coefficient rails, axle cornering stiffness is recovered to
\SI{3.2}{\percent} front and \SI{1.4}{\percent} rear and the peak factor to
\SI{11}{\percent}, against a railed peak factor and a
\SIrange{118}{132}{\percent} stiffness error for the inherited configuration. All results
are in simulation, one run per configuration.
\end{abstract}

\begin{IEEEkeywords}
Autonomous vehicles, system identification, parameter estimation, vehicle
dynamics, tires, predictive control, road vehicles, simulation.
\end{IEEEkeywords}

\IEEEpeerreviewmaketitle
